\part{Lattice Gauge Theory}

\section{Gauge fields on edges}

\begin{exercise}[Finite-group lattice gauge fields and gauge-invariant weights]
Let $G$ be a finite group and $X$ a finite oriented cell complex
whose faces have specified boundary edge words. A gauge field is
an assignment $g_e\in G$ with $g_{\bar e}=g_e^{-1}$.
The gauge transformations are the vertex actions of the section
TQFT, $g_e\mapsto h_{s(e)}g_eh_{t(e)}^{-1}$. The face holonomy
$U_f$ is the ordered product around its boundary, based at a vertex.
For a real class function $w$ define
\[
 Z_X(\beta_g)=|G|^{-|E|}\sum_g
       \exp\left(\beta_g\sum_f w(U_f)\right).
\]
\begin{enumerate}[label=(\alph*)]
\item Prove that changing the base vertex conjugates $U_f$ and
that the Gibbs measure is gauge invariant. State the additional
condition on $w$ for invariance under reversal of a face orientation.
\item For $G=\{\pm1\}$, $w(u)=u$, identify the Hamiltonian
$H=-J_g\sum_f U_f$ with $\beta_g=\beta J_g$.
On a connected graph fix a root and a spanning tree. Prove that
each field has a unique gauge transform trivial on the tree
if the transformation at the root is required to be the identity.
Determine the remaining root action.
\end{enumerate}
\end{exercise}

\section{Wilson loops}

\begin{exercise}[Wilson observables and the exact planar area law]
For a finite-dimensional representation $\rho:G\to\GL(V)$ and
an oriented closed edge path $C$ define the normalized Wilson loop
$W_\rho(C)=(\dim V)^{-1}\Tr\rho(U_C)$.
Use the defining one-dimensional representation for $G=\{\pm1\}$.
\begin{enumerate}[label=(\alph*)]
\item Prove gauge invariance and independence of the starting
vertex. For the $\Z/2$ model show that
$W(C)=\prod_{f\subset D}U_f$ when $C=\partial D$ is a simple
loop bounding a union of faces.
\item On a planar disc cellulation with free edge boundary
conditions show that the map from edge fields to face holonomies
is onto and that all fibres have size $2^{|E|-|F|}$.
One proof successively removes boundary faces; another uses
the tree elimination of the section Gauge fields on edges.
Deduce
\[
 Z_X=(\cosh\beta_g)^{|F|},\qquad
 \langle W(C)\rangle=(\tanh\beta_g)^{|D|}.
\]
\item Prove that any gauge-noninvariant single edge has zero
expectation, even when $\beta_g$ is large, by changing the gauge
at an endpoint. Compare this with the gauge-invariant loop.
\end{enumerate}
\end{exercise}

\section{Flat connections}

\begin{exercise}[Flat fields and the path groupoid]
A field is flat if $U_f=1$ for every face. The path groupoid
$\Pi(X)$ has vertices as objects and edge words as morphisms,
modulo immediate backtracking and face-boundary relations.
Regard $G$ as the one-object groupoid $BG_{\rm alg}$.
\begin{enumerate}[label=(\alph*)]
\item Prove that flat fields are functors
$\Pi(X)\to BG_{\rm alg}$ and that gauge transformations are
natural isomorphisms. Check that subdivision of an edge gives
an equivalent groupoid of flat fields.
\item On a connected complex with root $x$ prove that gauge
classes are conjugacy classes of homomorphisms
$\Aut_{\Pi(X)}(x)\to G$. Determine the automorphism group
of the field corresponding to $\rho$ as the centralizer
of $\im\rho$.
\item For a torus use its polygon with boundary word
$aba^{-1}b^{-1}$ to obtain commuting pairs $(A,B)\in G^2$
modulo simultaneous conjugation. For abelian $G$ count
$|G|^{2g}$ flat classes on an oriented genus-$g$ surface.
\end{enumerate}
\end{exercise}

\section{Dijkgraaf--Witten}

\begin{exercise}[Cocycle actions for finite gauge fields]
For the normalized three-cocycle $\omega$ of the section Pentagon
identities and a flat field on an ordered triangulated closed
oriented three-manifold define
\[
 Z_{G,\omega}(M)=|G|^{-n_0}\sum_{g\ {\rm flat}}
 \prod_{i<j<k<\ell}
 \omega(g_{ij},g_{jk},g_{k\ell})^{\epsilon_{ijk\ell}},
\]
where the product is over tetrahedra and $\epsilon$ compares the
vertex ordering with the orientation of $M$.
\begin{enumerate}[label=(\alph*)]
\item Recover the pointed state sum of the section Turaev--Viro.
Using the pentagon, check invariance under the $2$--$3$ move.
For the $1$--$4$ move show that each of the $|G|$ extensions
has the same total weight as the original tetrahedron.
\item Prove that replacing $\omega$ by $\omega\,\delta\nu$
leaves the closed sum unchanged: each oriented triangle factor
cancels with the factor from the tetrahedron on the other side.
For a manifold with boundary identify the uncancelled boundary
product explicitly.
\item For $\omega=1$ prove
$Z_{G,1}(M)=|\Hom(\pi_1(M),G)|/|G|$ when $M$ is connected,
where $\pi_1(M)=\Aut_{\Pi(M)}(x)$ for a triangulation.
Calculate $Z_{G,1}(S^3)$ and $Z_{G,1}(S^2\times S^1)$.
\end{enumerate}
\end{exercise}

\section{Classifying spaces}

\begin{exercise}[The nerve of a group and its classifying space]
The nerve of a category has as its $n$-simplices the strings of
$n$ composable arrows; face maps compose or omit end arrows and
degeneracy maps insert identities. Its geometric realization
glues a standard topological simplex to each such string by these
maps. Define $BG$ to be the realization of the nerve of
$BG_{\rm alg}$.
A homotopy between maps is a continuous map on the product with
$[0,1]$ restricting to the given maps at its ends.
\begin{enumerate}[label=(\alph*)]
\item Show that a flat field on an ordered simplicial complex
specifies a simplicial map to the nerve: an ordered simplex
$i_0<\cdots<i_n$ maps to
$(g_{i_0i_1},\dots,g_{i_{n-1}i_n})$.
Construct the homotopy induced by a vertex gauge transformation.
\item Let $\mathcal E G$ have object set $G$ and exactly one
arrow between each pair of objects. Prove its nerve is
contractible by inserting a fixed initial object. Prove that
the free $G$-action on its realization has quotient $BG$.
Deduce $\pi_1(BG)\simeq G$ and $\pi_n(BG)=0$ for $n>1$,
where $\pi_n$ denotes based maps from the $n$-sphere modulo
based homotopy.
\item For a connected finite cell complex build a map to $BG$
from a homomorphism of fundamental groups, first on its edges,
then across its faces and higher cells. Use the vanishing
higher homotopy groups to prove
$[X,BG]\simeq\Hom(\pi_1(X),G)/G$.
\end{enumerate}
\end{exercise}

\section{Higher gauge theory}

\begin{exercise}[Abelian two-form fields and gauge transformations between transformations]
For a finite cell complex $X$ let $C_n(X;\Z)$ be its free group of
oriented cells, with incidence boundary $\partial$, and put
$C^n(X;A)=\Hom(C_n(X;\Z),A)$ for an abelian group $A$.
The coboundary is $\delta c=c\partial$; define
$H^n(X;A)=\ker\delta/\im\delta$.
A flat two-form field is $b\in C^2(X;A)$ with $\delta b=0$.
A gauge transformation is $b\mapsto b+\delta a$, $a\in C^1(X;A)$,
and $a\mapsto a+\delta\lambda$, $\lambda\in C^0(X;A)$,
is a transformation between gauge transformations.
\begin{enumerate}[label=(\alph*)]
\item Verify $\partial^2=0$ for a simplicial boundary by pairing
the two orders of deleting vertices. Deduce $\delta^2=0$ and
that flat two-form fields modulo gauge transformations are
classified by $H^2(X;A)$.
\item For a two-cycle $S$ and a character $\chi:A\to U(1)$ define
$W_\chi(S)=\chi(b(S))$. Prove gauge invariance and invariance
under replacing $S$ by a homologous cycle when $b$ is flat.
\item Give the three-torus its product cell structure with one
vertex, three edges, three faces, and one three-cell.
Compute its cellular boundaries and prove
$H^2(T^3;A)=A^3$.
\end{enumerate}
\end{exercise}

\section{Homotopy types}

\begin{exercise}[Cocycle spaces as homotopy types of higher gauge fields]
For an abelian group $A$ and $n\ge1$ define a simplicial abelian
group by
$K(A,n)_r=Z^n(\Delta^r;A)$, the normalized simplicial
$n$-cocycles on the standard $r$-simplex, with restriction maps.
\begin{enumerate}[label=(\alph*)]
\item Prove that a simplicial map $X\to K(A,n)$ is the same
as an $n$-cocycle on $X$: restrict a cocycle to every simplex
and check the compatibility conditions.
\item Use the oriented prism on each simplex of $X\times[0,1]$
to show that homotopic maps give cohomologous cocycles.
Conversely extend a coboundary difference across the prism to
construct a homotopy. Obtain $[X,K(A,n)]=H^n(X;A)$ in
the simplicial homotopy category, allowing subdivision.
\item Apply this to spheres and derive
$\pi_j(K(A,n))=A$ for $j=n$ and zero otherwise.
Interpret the edge-field and two-form classifications as
$[X,K(A,1)]$ and $[X,K(A,2)]$ for abelian $A$.
\end{enumerate}
\end{exercise}

\begin{exercise}[Prediction: Wilson area decay and quantized higher flux sectors]
For the planar square $\Z/2$ gauge model with lattice spacing $a_0$
and $\beta_g=\beta J_g>0$, let $C$ enclose physical area
$\mathcal A=a_0^2|D|$. For the flat $\Z/2$ two-form model on a
spatial three-torus use the uniform superposition of every gauge
orbit as its zero-energy state, with constraints imposed by
commuting averages and flatness projections.
\begin{enumerate}[label=(\alph*)]
\item Derive
\[
 -\log\langle W(C)\rangle=\sigma_{\rm area}\mathcal A,\qquad
 \sigma_{\rm area}=-a_0^{-2}\log\tanh\beta_g.
\]
Obtain $\sigma_{\rm area}\sim2a_0^{-2}e^{-2\beta_g}$ at large
$\beta_g$ and calculate the ratio of expectations for loops
whose enclosed areas differ by one plaquette.
\item Derive eight ground states on the three-torus, labelled
by the three independent signs of the Wilson surfaces on
the coordinate two-tori. Construct the commuting surface
measurements distinguishing them using the pairing of the
section Higher gauge theory.
\item State which prediction requires free planar boundary
conditions and finite $\beta_g$, and which requires exact
flatness and periodic spatial boundaries. Identify loop and
surface expectation measurements in a lattice simulator that
test the two formulas.
\end{enumerate}
\end{exercise}
